\documentclass[t, aspectratio=169]{beamer} % Khai báo tỉ lệ cho Beamer biết
\usepackage[utf8]{vietnam}
\usepackage{amsthm,amsmath,amssymb}
\usepackage{lmodern}

% --- CẤU HÌNH LỀ (Phải đặt ở đây) ---
\def\slideContentLeftMargin{0.5cm}  % Lề trái
\def\slideContentRightMargin{1cm}   % Lề phải
\def\hustHeaderHeight{1cm}        % Lề trên (Né thanh tiêu đề)
\def\slideContentBotMargin{3.5cm}   % Lề dưới (Né logo/footer)


% Thay đổi [red] thành [blue] để đổi màu
\usepackage[blue, 169]{beamerthemeHUST} 

% Dùng dấu - thay cho tam giác ở các danh sách
\setbeamertemplate{itemize item}{--}
\setbeamertemplate{itemize subitem}{--}

% 1. Điều chỉnh vị trí cả khối title ở trang title
\renewcommand{\titlePageTitleX}{-0.5cm}
\renewcommand{\titlePageTitleY}{-0.5cm}

% 2. Căn lề TỪNG THÀNH PHẦN của trang title
\renewcommand{\hustTitleAlign}{\alignLeft} 
\renewcommand{\hustSubtitleAlign}{\alignLeft}
\renewcommand{\hustAuthorAlign}{\alignLeft}
\renewcommand{\hustInstAlign}{\alignLeft}


\title{BÁO CÁO TỔNG QUAN: FINANCIAL DATA PLATFORM (FDP)}
\author[Phạm Hải Nam]{Sinh viên thực hiện: Phạm Hải Nam - 20235791 \\ Giảng viên hướng dẫn: TS. Nguyễn Đức Toàn}  
\date{\today}

\begin{document}
	
	\hustintropage{} 
	\hustsectionpage{} 
	
	% Trang bìa
	\husttitlepage
	
	\begin{frame}[allowframebreaks]{Nội dung trình bày}
		\tableofcontents[hideallsubsections]
	\end{frame}
	
	\resetPageCount  
	\showPageNumber  
	
	\useStyleHustFull
	\useThinBar 
	
	% ==========================================
	\section{Phần I: Tổng quan đề tài}
	\begin{frame}
		\vfill
		\centering
		{\Huge \textbf{Phần I}} \\[1em]
		{\huge \textbf{Tổng quan đề tài}}
		\vfill
	\end{frame}
	
	\subsection{Đặt vấn đề}
	\begin{frame}{1. Đặt vấn đề}
		\textbf{Bối cảnh \& Nhu cầu:}
		\begin{itemize}
			\item Cần một hệ thống thu thập dữ liệu báo cáo tài chính công khai, chuẩn hóa thành các chỉ tiêu thống nhất.
			\item Yêu cầu kiểm tra chất lượng và lưu trữ dữ liệu có hỗ trợ truy vấn theo thời điểm lịch sử.
		\end{itemize}
		
		\vspace{1em}
		
		\textbf{Giải pháp - Hệ thống Financial Data Platform (FDP):}
		\begin{itemize}
			\item \textbf{Nguồn dữ liệu:} SEC Financial Statement Data Sets (doanh nghiệp niêm yết tại Mỹ).
			\item \textbf{Vai trò:} Đóng vai trò trung gian giữa nguồn dữ liệu công khai và người sử dụng thông qua giao diện web và API.
		\end{itemize}
	\end{frame}
	
	\subsection{Mục tiêu đề tài}
	\begin{frame}{2. Mục tiêu đề tài}
		Thiết kế và xây dựng hệ thống xử lý dữ liệu tài chính theo quy trình thống nhất:
		\begin{center}
			\textbf{Thu thập $\rightarrow$ Chuẩn hóa $\rightarrow$ Kiểm tra chất lượng $\rightarrow$ Lưu trữ $\rightarrow$ Truy vấn $\rightarrow$ Khai thác}
		\end{center}
		
		\vspace{0.5em}
		\textbf{Các mục tiêu cụ thể:}
		\begin{itemize}
			\item Thu thập dữ liệu tài chính từ SEC.
			\item Chuẩn hóa dữ liệu thô thành tập chỉ tiêu tài chính chuẩn hóa (canonical metrics).
			\item Kiểm tra và ghi nhận các vấn đề chất lượng dữ liệu.
			\item Lưu trữ dữ liệu lịch sử, duy trì các phiên bản (restatement).
			\item Khai thác: xem xu hướng, so sánh doanh nghiệp và lọc theo điều kiện.
			\item Cung cấp dữ liệu qua giao diện Web và API.
		\end{itemize}
	\end{frame}

	\subsection{Đối tượng sử dụng}
	\begin{frame}{3. Đối tượng sử dụng}
		Hệ thống hướng tới ba nhóm đối tượng chính:
		\begin{enumerate}
			\item \textbf{End User (Người dùng cuối):} Sử dụng giao diện web để tìm kiếm, tra cứu và xem thông tin tài chính của doanh nghiệp.
			\item \textbf{API Consumer (Bên tiêu thụ API):} Truy cập dữ liệu thông qua API phục vụ các mục đích xử lý hoặc tích hợp vào hệ thống khác.
			\item \textbf{System Operator (Người vận hành):} Phụ trách vận hành, theo dõi quá trình thu thập/xử lý dữ liệu, kiểm tra log và kiểm soát hệ thống.
		\end{enumerate}
	\end{frame}
	
	% ==========================================
	\useStyleLogoOnly 
	\useThickBar      
	
	\section{Phần II: Phân tích tập dữ liệu SEC}
	\begin{frame}
		\vfill
		\centering
		{\Huge \textbf{Phần II}} \\[1em]
		{\huge \textbf{Phân tích tập dữ liệu SEC}}
		\vfill
	\end{frame}
	
	\subsection{Phân tích file sub.txt và num.txt}
	\begin{frame}{Phân tích file sub.txt và num.txt}
		\begin{columns}[T]
			\begin{column}{0.48\textwidth}
				\begin{block}{sub.txt (Thông tin báo cáo)}
					\begin{itemize}
						\item \textbf{CIK:} Mã nhận diện công ty từ SEC.
						\item \textbf{adsh:} Khóa chính, mã truy cập hồ sơ.
						\item \textbf{sic:} Mã ngành nghề kinh doanh.
						\item \textbf{form:} Loại báo cáo (10-K, 10-Q).
						\item \textbf{fy \& fp:} Năm và kỳ tài chính.
					\end{itemize}
				\end{block}
			\end{column}
			
			\begin{column}{0.48\textwidth}
				\begin{block}{num.txt (Dữ liệu số liệu)}
					\begin{itemize}
						\item \textbf{adsh:} Khóa ngoại liên kết sub.txt.
						\item \textbf{tag:} Tên chỉ số (NetIncomeLoss,...).
						\item \textbf{ddate:} Ngày chốt sổ.
						\item \textbf{qtrs:} Thời gian phát sinh (0: 1 thời điểm, 1: 1 quý, 4: 1 năm).
						\item \textbf{value:} Giá trị thực tế.
					\end{itemize}
				\end{block}
			\end{column}
		\end{columns}
	\end{frame}

	\subsection{Phân tích file tag.txt và pre.txt}
	\begin{frame}{Phân tích file tag.txt và pre.txt}
		\begin{columns}[T]
			\begin{column}{0.48\textwidth}
				\begin{block}{tag.txt (Danh mục tham chiếu)}
					\begin{itemize}
						\item Định nghĩa cấu trúc cho chỉ số.
						\item \textbf{datatype:} Kiểu dữ liệu (monetary, shares, pure).
						\item \textbf{crdr:} Cờ quy đổi dấu C/D (Credit/Debit).
						\item \textbf{tlabel \& doc:} Nhãn hiển thị UI và mô tả nghiệp vụ.
					\end{itemize}
				\end{block}
			\end{column}
			
			\begin{column}{0.48\textwidth}
				\begin{block}{pre.txt (Cấu trúc trình bày)}
					\begin{itemize}
						\item Hỗ trợ tái tạo giao diện gốc.
						\item \textbf{stmt:} Bảng chứa số liệu (BS, IS, CF).
						\item \textbf{line:} Thứ tự dòng.
						\item \textbf{inpth:} Độ thụt lề (cấp bậc khoản mục).
						\item \textbf{negating:} Đổi dấu khi hiển thị UI (1: đổi).
					\end{itemize}
				\end{block}
			\end{column}
		\end{columns}
	\end{frame}
	
	% ==========================================
	\useStyleHustFull
	\useThinBar 
	
	\section{Phần III: Nghiên cứu Case Study}
	\begin{frame}
		\vfill
		\centering
		{\Huge \textbf{Phần III}} \\[1em]
		{\huge \textbf{Nghiên cứu Case Study}}
		\vfill
	\end{frame}
	
	\subsection{Các nền tảng dữ liệu tham khảo}
	\begin{frame}{Các nền tảng dữ liệu tham khảo (1)}
		\textbf{1. Nasdaq Data Fabric}
		\begin{itemize}
			\item \textbf{Mô hình:} "Managed pipeline" - quản lý nạp dữ liệu, chất lượng và hạ tầng.
			\item \textbf{Đặc điểm:} Cung cấp catalog tập trung, phân quyền chi tiết qua một API duy nhất.
			\item[$\rightarrow$] Ứng dụng data catalog tập trung và phân quyền cơ bản.
		\end{itemize}
		
		\vspace{1em}
		
		\textbf{2. Polygon.io}
		\begin{itemize}
			\item \textbf{Dịch vụ:} Cung cấp dữ liệu thị trường (stock, crypto, forex).
			\item \textbf{Đặc điểm:} API REST versioned rất chuẩn mực (\texttt{/v1/\{resource\}/...}) và WebSocket.
			\item[$\rightarrow$] Cấu trúc REST endpoint chuẩn để thiết kế API cho FDP.
		\end{itemize}
	\end{frame}

	\begin{frame}{Các nền tảng dữ liệu tham khảo (2)}
		\textbf{3. Economatica}
		\begin{itemize}
			\item \textbf{Mô hình:} Nền tảng phân tích tài chính sâu (40 năm dữ liệu).
			\item \textbf{Tính năng:} Screener (lọc đa điều kiện), Comparative matrix (so sánh dữ liệu).
			\item[$\rightarrow$] Gợi ý tính năng query lọc và vẽ biểu đồ so sánh.
		\end{itemize}
		
		\vspace{1em}
		
		\textbf{4. AlphaSense}
		\begin{itemize}
			\item \textbf{Mô hình:} Tìm kiếm bằng AI/NLP trên tài liệu phi cấu trúc (transcript, news).
			\item[$\rightarrow$] Khác nhóm bài toán nhưng gợi ý hướng mở rộng AI/Semantic Search trong tương lai.
		\end{itemize}
	\end{frame}
	
	% ==========================================
	\useStyleLogoOnly 
	\useThickBar      
	
	\section{Phần IV: Phạm vi của đề tài}
	\begin{frame}
		\vfill
		\centering
		{\Huge \textbf{Phần IV}} \\[1em]
		{\huge \textbf{Phạm vi của đề tài}}
		\vfill
	\end{frame}
	
	\subsection{Phạm vi chức năng}
	\begin{frame}{1. Phạm vi chức năng}
		\begin{itemize}
			\item \textbf{Thu thập (Ingestion):} Dùng khóa (CIK + adsh), ghi log chi tiết.
			\item \textbf{Chuẩn hóa:} Ánh xạ về canonical metric theo mức ưu tiên; xử lý theo thời đoạn (qtrs) và dấu (crdr).
			\item \textbf{Kiểm tra chất lượng:} Dùng khóa (cik + tag + ddate + qtrs + adsh) phát hiện trùng, gắn cờ lỗi (không xóa).
			\item \textbf{Lưu trữ lịch sử:} Báo cáo điều chỉnh lưu song song, không ghi đè; hỗ trợ versioning.
			\item \textbf{Truy vấn \& Phân tích:} Tìm kiếm DN, xem lịch sử/xu hướng, so sánh (2-5 DN), tính chỉ số dẫn xuất.
			\item \textbf{Kênh cung cấp:} Giao diện Web (Desktop) và REST API.
		\end{itemize}
	\end{frame}

	\subsection{Phạm vi dữ liệu}
	\begin{frame}{2. Phạm vi dữ liệu}
		\begin{block}{Quy mô tham chiếu (SEC Data)}
			\begin{itemize}
				\item \textbf{Quy mô doanh nghiệp:} 30 doanh nghiệp (đa dạng mã ngành - sic).
				\item \textbf{Độ sâu thời gian:} Tối thiểu 12 quý dữ liệu lịch sử.
				\item \textbf{Dữ liệu chỉ tiêu:} Tập canonical metrics xác định trước.
				\item \textbf{Dữ liệu thị trường:} Bổ sung dữ liệu lịch sử giá cổ phiếu.
			\end{itemize}
		\end{block}
		
		\vspace{1em}
		
		\begin{alertblock}{Mục tiêu chất lượng dữ liệu}
			Đạt tối thiểu \textbf{70\% coverage} (độ phủ) đối với ít nhất \textbf{6 canonical metrics}.
		\end{alertblock}
	\end{frame}
	
	\hustthankyou
	
\end{document}
